\section{Implementation}
\label{sec:implementation}

HStore is written in Java~25 as three modules. The storage engine (about 12,000 non-blank lines) depends only on the Java base module. The database layer (about 7,900 lines) adds schemas and typed properties, property indexes, tenants and users, evidence and valid-time queries, a three-way branch merge, an embedding index fed from the change feed, higher-order operators over incidence sets, and HQL, a query language with a cost-based planner. The server (about 2,600 lines) provides a command-line tool, a line-oriented TCP protocol and a web console. The same code runs on the JVM and as a GraalVM native image, which is how the server ships in its container image.

The engine leans on recent additions to the Java platform. Segment maps and the page directory use the foreign memory API; each segment map has its own shared arena and is closed as soon as the segment is remapped or deleted, so the disk space of a deleted segment comes back at once instead of when the garbage collector finds the mapping. A per-query I/O meter, which enforces a page budget per query, is bound with scoped values rather than thread-locals. Server connections and feed subscribers run on virtual threads. Commits can therefore be issued from thousands of threads, and the group committer runs its sync on whichever committing thread leads the batch.

A Java thread that is interrupted during a read or write on a \code{FileChannel} closes the channel for every thread. Before this was handled, one request that timed out during a commit closed the log and every later commit failed. All shared files are therefore accessed through a wrapper that clears the interrupt flag for the length of the call, restores it afterwards, and reopens the file and retries if the channel was closed anyway; every read and write is positional, so a retry touches the same bytes.

\paragraph{Testing.} The engine's tests inject a crash at each of fourteen named points: the seven steps of a commit from the first page write to publication, two steps of compaction, four of a checkpoint and the force of the log. Each commit crash point is run with both log modes; the test asserts that the reopened store holds either the state before the commit or the state after it, never a mix, and that the forward and reverse indexes agree. A second set of tests models a power failure instead of a process crash. Every file the engine writes is opened through one wrapper, which lets a test record each write, each sync of a file or a memory map and each sync of a directory, cut the power at any of these events, and rebuild the directory as a disk could have kept it: only what was synced, the last unsynced write torn in half or with a flipped byte, or a random mix, with files whose directory was never synced possibly missing. The tests then check that the recovered commits are a prefix of the attempted ones, that every commit acknowledged under synchronous durability survived, that no aborted work appears and that the change feed holds exactly one event per recovered commit. They found five bugs, all of them in files outside the log: a format marker written without a sync, a segment directory that was never synced after a file was created in it, a page directory header that could keep an unsynced first page number, and two ways in which zeroed or damaged change-feed frames were read as valid. After random workloads, the suites replay the log against the generations in memory and compare them entry by entry, and a verifier checks every reachable page (checksums, key order, balance, summaries and reference counts), which the \code{hstore check} command also runs. Finally, we kill the native server with SIGKILL during a bulk load, reopen it and check the counts.
